% BHU PYQ -- Professional Exam Template
\documentclass[11pt,a4paper]{article}

\usepackage[a4paper,top=2.5cm,bottom=2.5cm,left=2.8cm,right=2.8cm,headheight=14pt]{geometry}
\usepackage{amsmath,amssymb}
\usepackage[T1]{fontenc}
\usepackage[utf8]{inputenc}
\usepackage{lmodern}
\usepackage{microtype}
\usepackage{fancyhdr}
\usepackage{enumitem}
\usepackage{listings}
\usepackage{hyperref}
\usepackage{lastpage}
\usepackage{parskip}

\hypersetup{colorlinks=true,urlcolor=black,linkcolor=black,
  pdftitle={BVCA-402: Probability and Statistics -- BHU PYQ}}

\pagestyle{fancy}
\fancyhf{}
\renewcommand{\headrulewidth}{0.4pt}
\renewcommand{\footrulewidth}{0.4pt}
\fancyhead[L]{\small   \textit{Banaras Hindu University}}
\fancyhead[R]{\small   \textit{BVCA-402: Probability and Statistics}}
\fancyfoot[C]{\small Page  \thepage\ of \pageref{LastPage}}

\lstset{
  basicstyle=\small tfamily,
  frame=single,
  breaklines=true,
  columns=flexible,
  keepspaces=true
}


\newlist{parts}{enumerate}{2}
\setlist[parts,1]{label=(\alph*),leftmargin=*,itemsep=4pt,topsep=3pt}
\setlist[parts,2]{label=(\roman*),leftmargin=*,itemsep=2pt,topsep=2pt}

\newcommand{\pts}[1]{\hfill   \textit{[#1]}}


\begin{document}

\begin{center}
  {\large\bfseries BANARAS HINDU UNIVERSITY}\\[2pt]
  {\normalsize B.Voc. Semester IV Examination, 2023-24}\\[6pt]
  \rule{0.8\linewidth}{0.5pt}\\[6pt]
  {\large\bfseries Paper: BVCA-402 --- Probability and Statistics}\\[4pt]
  {\normalsize Time: Three Hours \hfill Maximum Marks: 70}
\end{center}

\rule{\linewidth}{0.4pt}

\smallskip



\noindent   \textit{Instructions: Answer the questions in each section as instructed. All questions carry marks as indicated.}

\bigskip




\noindent {\large\bfseries SECTION --- A}
\rule{\linewidth}{0.4pt}\smallskip

Long-Answer Questions. Answer any two questions: \hfill   \textbf{2 $ \times$ 12 = 24 Marks}

\medskip



\noindent   \textbf{Question 1.} Derive the mean and variance of the Normal distribution $N(\mu, \sigma^2)$.

\medskip



\noindent   \textbf{Question 2.} What is a simple linear regression model? Obtain the ordinary least squares expressions for the estimators $eta_0$ and $eta_1$.

\medskip



\noindent   \textbf{Question 3.} What is a cumulative distribution function? Explain the properties of the distribution function with their mathematical proofs.

\medskip



\noindent   \textbf{Question 4.} What is a likelihood function? Explain the theory of how to obtain the maximum likelihood estimator (MLE) of the parameter of any distribution. In random sampling from a Normal distribution $N(\mu, \sigma^2)$, find the maximum likelihood estimator (MLE) for:

\begin{parts}
  \item $\mu$ when $\sigma^2$ is known,
  \item $\sigma^2$ when $\mu$ is known.
\end{parts}

\bigskip


\noindent {\large\bfseries SECTION --- B}
\rule{\linewidth}{0.4pt}\smallskip




\noindent   \textbf{Question 5.} Attempt any six parts (Answer in about 200 words each): \hfill   \textbf{6 $ \times$ 6 = 36 Marks}

\begin{parts}
  \item Calculate the mean and standard deviation of the continuous random variable $X$ if its probability density function $f(x)$ is given by:
  \[ f(x) = 
\begin{cases} \frac{3}{56}x^2 &  \text{if } 2 < x < 4 \\ 0 &  \text{otherwise} \end{cases} \]
  \item From a city population, the probability of selecting:
  
\begin{parts}
    \item a male or a smoker is $7/10$,
    \item a male smoker is $2/5$,
    \item a male, if a smoker is already selected, is $2/3$.
  \end{parts}
  Find the probability of selecting:
  
\begin{parts}
    \item a non-smoker,
    \item a male,
    \item a smoker, if a male is first selected.
  \end{parts}
  \item The contents of three urns I, II, and III are as follows:
  
\begin{itemize}[label=--,leftmargin=1.5em,itemsep=2pt]
    \item Urn I: 1 white, 2 black, and 3 red balls
    \item Urn II: 2 white, 1 black, and 1 red balls
    \item Urn III: 4 white, 5 black, and 3 red balls
  \end{itemize}
  One urn is chosen at random and two balls are drawn from it. They happen to be white and red. What is the probability that they come from:
  
\begin{parts}
    \item Urn I,
    \item Urn II,
    \item Urn III?
  \end{parts}
  \item Define the chi-square test statistic. What are the essential conditions for the validity of the chi-square test?
  \item What are non-parametric tests? Define the Kolmogorov-Smirnov statistic.
  \item For any three events $A$, $B$, and $C$, prove the following:
  
\begin{parts}
    \item $P(A \cup B \mid C) = P(A \mid C) + P(B \mid C) - P(A \cap B \mid C)$
    \item $P(B \mid A) \le P(C \mid A)$ provided that $B \subseteq C$ and $P(A) > 0$.
  \end{parts}
  \item Discuss statistical hypothesis. Define null and alternative hypothesis with an illustrative example.
  \item State and define Chebyshev's inequality and Markov's inequality.
  \item Define Statistics. According to you, where can it be applied? Give at least three practical examples.
\end{parts}

\bigskip


\noindent {\large\bfseries SECTION --- C}
\rule{\linewidth}{0.4pt}\smallskip




\noindent   \textbf{Question 6.} Attempt all parts: \hfill   \textbf{10 $ \times$ 1 = 10 Marks}

\begin{parts}
  \item In an experiment, a coin is tossed five times. Write down the sample space.
  \item State the additive property of the Poisson distribution.
  \item Define the Weak Law of Large Numbers (WLLN).
  \item Define Type-II error in hypothesis testing.
  \item Who discovered the Normal distribution?
  \item Write down the probability density function (PDF) of a standard normal variate.
  \item Give two practical examples of a Poisson distribution.
  \item For any three events $A$, $B$, and $C$, write the expression for $P(A \cup B \cup C)$.
  \item Define exhaustive events.
  \item Write down the mean and variance of a standard normal variate.
\end{parts}

\vfill
\rule{\linewidth}{0.4pt}

\begin{center}\small
\href{https://bhu-exam.vercel.app/}{Click here to visit our website}\\[4pt]\href{https://me-aryan.vercel.app/}{Click here to visit our contributor}\\[4pt]\href{https://quashan-me.vercel.app/}{Click here to visit our contributor}
\end{center}
\end{document}
